\documentclass{article}

\usepackage[preprint]{neurips_2026}

\usepackage[utf8]{inputenc}
\usepackage[T1]{fontenc}
\usepackage{hyperref}
\usepackage{url}
\usepackage{booktabs}
\usepackage{amsmath}
\usepackage{amsfonts}
\usepackage{nicefrac}
\usepackage{microtype}
\usepackage{xcolor}
\usepackage{graphicx}

\newcommand{\good}[1]{\textcolor{green!40!black}{#1}}
\newcommand{\bad}[1]{\textcolor{red!40!black}{#1}}

\title{Evaluating Author-Level Citation Metrics on a Closed Citation Graph:
A Pre-registered Comparison}

\author{%
  Anonymous Author(s) \\
  CS6101 Project \\
}

\begin{document}

\maketitle

\begin{abstract}
Ranking researchers by citation impact shapes hiring, funding, and tenure decisions, yet the
sensitivity of such rankings to the choice of metric is rarely measured in a pre-registered
framework on a single closed dataset. We evaluate seven author-level metrics, from field-share
sums and per-paper averages to citation deflation and core-based indices, against a held-out
future-citation outcome on a closed graph of 14,462 Information Retrieval (IR) papers from five
venues (2000--2022). Eight hypotheses were fixed in writing before any metric was computed.
Dividing the field-share sum (A1) by paper count (M1) costs $0.083$ in rank accuracy
($\rho$: $0.708 \to 0.625$; bootstrap 95\% CI of the difference $[-0.099, -0.067]$) and leaves the
correlation between score and cohort birth year unchanged (difference CI $[-0.024, +0.010]$),
reproducing the documented behaviour of the mean normalised citation score. A Petersen-style
citation deflator ($t_0 = 2021$) costs $0.014$ in accuracy and reorders almost nothing (top-100
overlap 97 of 100 on OpenAlex and 96 of 100 on Semantic Scholar, first nine ranks identical). The top-3 mean (M3) is the closest non-trivial
alternative to A1 ($\rho = 0.692$) but measures best-paper concentration rather than
contribution. Each of these findings needs a further test: the age correlation is confounded by a
mechanical relation between a paper's age and how much future-citation window it can reach, and
no author-level metric can remove that confound. We release the pre-registrations, code, and
per-author scores so that the comparison can be repeated on other fields.
\end{abstract}

\section{Introduction}

Ranking researchers within a scientific field matters for hiring committees, funding agencies, and
tenure panels. The dominant family of solutions normalises raw citation counts by field and
publication year, producing indicators such as the field-weighted citation impact (FWCI) or the
mean normalised citation score (MNCS) \citep{waltman2016review, waltman2019chapter}. Alternative
families---percentile-based indicators \citep{bornmann2011further, bornmann2013percentiles}, core-based indices
\citep{hirsch2005index, egghe2006gindex}, and network-centrality scores \citep{ding2011pagerank,
west2010eigenfactor}---are well theorised but seldom compared head-to-head on one dataset against one
outcome.

Citation counts grow with paper age, so any ranking rewards earlier cohorts, and the usual remedy
is to normalise by field and year \citep{waltman2019chapter}, and academic-age normalisation, which
divides by years since first publication, is now offered as a platform-level service
\citep{anonymous2026platform}. \citet{petersen2019inflation} proposed a further correction, a \emph{citation
deflator} that divides each citation by the reference volume of its citing year, in analogy to a
currency deflator. Whether these corrections change who is ranked
where, and at what cost in accuracy, is an empirical question that depends on the field and on
the outcome.

This paper reports such an evaluation on one closed graph. Eight hypotheses were fixed in writing
before the corresponding metrics were computed, every verdict is reported whatever its direction,
and all point estimates come with bootstrap confidence intervals from a common author resample so
that metric-to-metric differences are paired. The outcome is a held-out window (2018--2022) that
no metric can see. A 20-run noise floor (20\% of edges dropped) and an independent Semantic
Scholar mirror of the same papers test robustness.

The scope is deliberately narrow. This is a \emph{measurement comparison}, not a field-wide
validation: it uses one field, one birth window, one citation source at a time, and a single
5-year outcome. Its contribution is to show which of the standard corrections are distinguishable
from noise at this scale, and---more importantly for practice---which trade-offs the usual
normalisations do and do not buy.

\section{Related Work}

\subsection{Normalisation families}

\citet{waltman2016review} and \citet{waltman2019chapter} survey field-normalised citation indicators. Two
families dominate. \emph{Ratio-based} indicators divide a paper's citations by the mean of its
field-year cohort and aggregate either by summing (size-dependent) or averaging
(size-independent: MNCS, FWCI). Our A1 is a size-dependent member of this family with an
author-level denominator. The \emph{percentile-based} family
\citep{tijssen2002benchmark, vanleeuwen2003holy, leydesdorff2011turning, albarran2011skewness} values a paper by its position
in its cohort distribution and is less sensitive to fat tails. Ties are known to be a problem
there; \citet{waltman2013percentile} compare tie-handling schemes. Percentile scoring is not implemented in
this study and is the first item of future work.

\subsection{Citation inflation}

\citet{petersen2019inflation} report that the reference volume produced per year grows at roughly 5.6\%
annually and argue that this ``citation inflation'' understates the impact of older generations.
Their correction is $D(t) = R(t_0)/R(t)$, with $R(t)$ the reference supply in year $t$. The growth
kinetics are analysed by \citet{pan2018memory} and the case for deflation as a complement to
year-normalisation is made by \citet{ioannidis2014most}. Section~\ref{sec:deflator} describes the
variant of this correction used here.

\subsection{Productivity normalisation}

Dividing an author's citations by their paper count yields a per-paper average, which is the
author-level analogue of MNCS. Its failure modes are known: instability at small $n$, domination
by single highly cited papers, and rank reversals against total-influence measures
\citep{waltman2019chapter}.

\subsection{Core-based indices}

The $h$-index \citep{hirsch2005index} and its variants---the $g$-index \citep{egghe2006gindex}, $w$-index
\citep{wu2010windex}, and $m$-quotient \citep{burrell2007core}---are widely used. \citet{alonso2009review} and
\citet{wildgaard2014108} review their properties. The $h$-index grows with career length.

\subsection{Credit allocation and other signals}

\citet{vavrycuk2018influential} argues that author position carries information about contribution, which
matters for IR venues where advisor--student papers are common. \citet{ding2011pagerank} applied
PageRank to an author co-citation network of 108 IR authors, the closest published precedent for
the present scale. The \emph{CD index} \citep{funk2017cdindex, wu2019teams} measures disruption rather than
influence and is orthogonal to citation counts.

\section{Method}

\subsection{Data}

\paragraph{Corpus.} 14,462 papers from SIGIR, CIKM, WSDM, TOIS, and IRJ, published 2000--2022,
with author names in listed order and 345,866 distinct references in the source records. Citation
edges are year-stamped.

\paragraph{Citation sources.} Two graphs over the same papers are used:
\begin{itemize}
  \item \textbf{OpenAlex} (primary): 54,351 in-corpus citation edges, 1,430 authors after filtering.
  \item \textbf{Semantic Scholar} (cross-check): 13,548 edges, 1,456 authors.
\end{itemize}
The Semantic Scholar graph is much thinner, with tie fractions up to 100\% (Table~\ref{tab:s2}), so
it is used to test the \emph{direction} of an effect, not its magnitude.

\paragraph{Scoring and outcome.} A paper born in year $y$ is scored on citations received in
$[y, y+\Delta]$ with $\Delta = 3$. Authors need at least 3 scored papers born in 2006--2014. The
outcome is the citations each author's papers (born $\leq 2017$) received in 2018--2022,
self-citations removed (a citation from a paper whose first author is the author themselves) and
credit split $1/n$ across co-authors. The outcome window lies entirely after the latest scoring
window (2014--2017), so no scored citation is counted as outcome.

\subsection{Metrics}

Table~\ref{tab:defs} defines the metrics. All return (author, score, $n_{\text{papers}}$,
raw citations) and share the 3-paper minimum, so rows differ only in aggregation.

\begin{table}[h]
\caption{Metric definitions. $c(p)$ = citations to paper $p$ inside its scoring window,
  $N_F(W)$ = total scored citations to all papers whose scoring window is $W$,
  $n_a(p)$ = number of authors of $p$, and $\bar{s}$ = field median paper share.}
\label{tab:defs}
\centering
\small
\begin{tabular}{ll}
\toprule
\textbf{Metric} & \textbf{Definition} \\
\midrule
A1 & $\sum_p \dfrac{c(p)}{N_F([y_p, y_p+\Delta])\, n_a(p)}$ \\[8pt]
M1 & A1 $/\, n_{\text{papers}}$ \\[4pt]
M2 & median over $p$ of the paper share \\[4pt]
M3 & mean of the author's 3 largest paper shares \\[4pt]
M4 & $(n \cdot \text{M1} + k\bar{s})/(n + k)$, $k = 3$ \\[4pt]
B1 & $\sum_p c(p)/n_a(p) \big/ n_{\text{papers}}$ \\[4pt]
B2 & number of papers with $\geq 3$ scored citations \\
\midrule
D1 & A1 with each citation weighted by $D(y_{\text{citing}})$ \\
D2 & D1 $/\, n_{\text{papers}}$ \\
B1$'$ & B1 with deflated citations \\
\bottomrule
\end{tabular}
\end{table}

\subsection{Citation deflator}
\label{sec:deflator}

We follow \citet{petersen2019inflation} with
\begin{equation}
  D(t) = \frac{R(t_0)}{R(t)}, \qquad
  R(t) = \frac{\text{outside-pool references emitted by pool papers of year } t}
               {\text{number of pool papers of year } t},
\end{equation}
restricting references to those citing works outside the pool, so that $R(t)$ does not measure
the growth of our own corpus.

Amendment~1 was recorded before any deflated metric was computed. $t_0 = 2021$, the last year of
the outcome window, so outcome-window citations carry weight near 1; years before 2005 are
excluded because 41--54\% of 2001--2003 papers have truncated reference lists
(\texttt{ir\_deflator\_truncation.csv}); and 0 of 10,896 in-window OpenAlex edges (0 of 577 on
Semantic Scholar) fall in an undefined year, so no author loses evidence. Two further variants
are reported as cross-checks: $D_{\text{naive}}$, built on raw reference counts without the
per-paper normalisation, and $D_{\text{distinct}}$, built on distinct outside works cited. As a
guard on the implementation, setting $D(t) \equiv 1$ must reproduce A1 exactly; it does, at
$\rho = 1.00000000$ with maximum relative error $0.00\text{e}0$.

\subsection{Evaluation and noise floor}

For each metric we report the Spearman correlation with the outcome ($\rho_{\text{fut}}$), the
correlation between score and the author's mean scored-paper year ($\text{corr}_{\text{yr}}$;
negative means earlier cohorts score higher), mean outcome of the top 10 and top 100, the tie
fraction, and agreement with A1 (Spearman correlation and top-100 overlap). Confidence intervals
come from 2000 bootstrap resamples of the 1,430 authors, drawn once and shared across metrics so
that differences against A1 are paired; intervals are percentile intervals.

The noise floor drops 20\% of citation edges at random in 20 runs. A1's minimum correlation with
the full ranking is 0.986 and its minimum top-100 overlap is 93 (Table~\ref{tab:noise}), far
tighter than the $|\Delta\rho| < 0.03$ band fixed in the plan, so every between-metric difference
is judged against that floor.

\subsection{Pre-registered hypotheses}

All hypotheses were written in \texttt{METRIC\_PLAN.md} and \texttt{METRIC\_PLAN\_DEFLATOR.md}
before the corresponding metrics were computed, and the deflator protocol was amended once
(Amendment~1, fixing $t_0$ and the truncation rule) before any deflated metric was computed. The
verdict on each hypothesis below is the one recorded in \texttt{RESULTS.md} and
\texttt{RESULTS\_DEFLATOR.md} at the time it was tested; where this paper disagrees with a
recorded verdict (H8) the disagreement is stated in place.
\begin{description}
  \item[H1.] Per-paper normalisation (M1) reduces $|\text{corr}_{\text{yr}}|$ relative to A1.
  \item[H2.] Per-paper normalisation lowers $\rho_{\text{fut}}$ below A1's 0.708.
  \item[H3.] Shrinkage toward the field mean (M4, $k=3$) recovers most of H2's loss.
  \item[H4.] The top-3 mean (M3) keeps the highest top-10 outcome among per-paper metrics, and a
    core-count index (B2) has the worst tie fraction.
  \item[H5.] The deflator reduces $|\text{corr}_{\text{yr}}|$ relative to A1.
  \item[H6.] The deflator costs less accuracy than M1 did, i.e.\ $\rho_{\text{fut}}(\text{D1}) >
    0.625$.
  \item[H7.] The deflator reorders substantially: $\rho(\text{D1}, \text{A1}) \in [0.80, 0.95]$
    and top-100 overlap below 70.
  \item[H8.] The field-share normalisation adds information beyond deflation: D1 beats B1$'$ on
    $\text{corr}_{\text{yr}}$, i.e.\ D1's value is the one closer to zero. (\texttt{RESULTS\_DEFLATOR.md}
    records this hypothesis with the opposite inequality, ``D1 has a more negative
    $\text{corr}_{\text{yr}}$''; that version cannot fail against the observed numbers and is not
    the one tested here. See the H8 paragraph.)
\end{description}

\section{Results}

\subsection{Productivity normalisation (H1--H4)}

\begin{table}[h]
\caption{OpenAlex (54,351 edges, 1,430 authors). $\text{corr}_{\text{yr}}$: negative = earlier
  cohorts score higher.}
\label{tab:oa}
\centering
\footnotesize
\setlength{\tabcolsep}{3.5pt}
\begin{tabular}{lrrrrr}
\toprule
\textbf{Metric} & $\rho_{\text{fut}}$ & $\text{corr}_{\text{yr}}$ & \textbf{Top-10} & \textbf{Top-100} & \textbf{Ties} \\
\midrule
A1 & 0.708 & $-$0.178 & 96.5 & 30.5 & 22\% \\
M1 & 0.625 & $-$0.185 & 22.3 & 20.4 & 19\% \\
M2 & 0.513 & $-$0.071 & 7.7  & 13.8 & 82\% \\
M3 & 0.692 & $-$0.183 & 46.3 & 27.8 & 33\% \\
M4 & 0.302 & $-$0.139 & 26.4 & 20.4 & 17\% \\
B1 & 0.646 & $-$0.080 & 19.8 & 22.3 & 48\% \\
B2 & 0.689 & $-$0.071 & 85.1 & 32.0 & 98\% \\
\bottomrule
\end{tabular}
\end{table}

\begin{table}[h]
\caption{Bootstrap 95\% CIs, OpenAlex. $d_\rho$ = paired difference in $\rho_{\text{fut}}$ from
  A1; $d_{\text{yr}}$ = paired difference in $\text{corr}_{\text{yr}}$ from A1.}
\label{tab:ci}
\centering
\footnotesize
\setlength{\tabcolsep}{3pt}
\begin{tabular}{lrrr}
\toprule
\textbf{Metric} & $d_\rho$ \textbf{CI} & $d_{\text{yr}}$ \textbf{CI} & \textbf{Overlap} \\
\midrule
M1 & [$-$0.099, $-$0.067] & [$-$0.024, $+$0.010] & 48 \\
M2 & [$-$0.226, $-$0.163] & [$+$0.071, $+$0.145] & 31 \\
M3 & [$-$0.027, $-$0.006] & [$-$0.017, $+$0.009] & 72 \\
M4 & [$-$0.455, $-$0.357] & [$-$0.008, $+$0.086] & 48 \\
B1 & [$-$0.079, $-$0.046] & [$+$0.079, $+$0.116] & 47 \\
B2 & [$-$0.036, $-$0.001] & [$+$0.081, $+$0.134] & 79 \\
\bottomrule
\end{tabular}
\end{table}

\paragraph{H1: failed.} $\text{corr}_{\text{yr}}$ moves from $-0.178$ (A1) to $-0.185$ (M1), the
wrong direction, and the paired CI for the change, $[-0.024, +0.010]$, contains zero. On Semantic
Scholar it moves from $-0.002$ to $+0.010$ (paired CI $[+0.004, +0.021]$), the same sign of change
on a graph where both values are within $\pm 0.01$ of zero. M1
does remove volume: the correlation between score and paper count falls from 0.61 to 0.33. It
removes nothing from the score's relation to birth year.

\paragraph{H2: held.} $\rho_{\text{fut}}$ falls from 0.708 to 0.625, a difference of $0.083$
(paired CI $[-0.099, -0.067]$), eight times A1's bootstrap half-width. The top-10 mean falls from
96.5 to 22.3, and M1 shares 48 of A1's top 100. M1 correlates with B1, plain per-paper citations, at
$\rho = 0.923$ and shares 47 of 100 with it, so averaging and dropping the field-share
denominator move the ranking in the same broad direction but are not interchangeable.

\paragraph{H3: failed.} M4 gives $\rho_{\text{fut}} = 0.302$, the worst value in
Table~\ref{tab:oa}, with a difference CI of $[-0.455, -0.357]$ against A1. With $k=3$ and a
3-paper minimum, shrinkage moves the score of an author with the minimum three papers by one half
and that of an author with six papers by one third, pulling the ranking toward mid-size records
without adding accuracy. Only $k=3$ was tested.

\paragraph{H4: half held.} M3 gives $\rho_{\text{fut}} = 0.692$ (difference CI
$[-0.027, -0.006]$) with the highest top-10 mean among the metrics other than A1 and B2 (46.3); it is
essentially A1 with the tail cut off ($\rho = 0.969$ with A1, 72 of 100 shared). It does not beat
A1. B2 has the worst ties (98\%) as predicted, but its accuracy prediction held only on Semantic
Scholar ($\rho = 0.223$); on OpenAlex it reaches 0.689, above M1 and B1.

\subsection{Noise floor}

\begin{table}[h]
\caption{Noise floor: 20 runs with 20\% of edges dropped. $\rho_{\min}$ and overlap minima are
  with respect to the full-data ranking.}
\label{tab:noise}
\centering
\footnotesize
\setlength{\tabcolsep}{4pt}
\begin{tabular}{lrrr}
\toprule
\textbf{Metric} & $\rho_{\min}$ & $\rho$ median & Overlap min \\
\midrule
A1 & 0.986 & 0.988 & 93 \\
M1 & 0.980 & 0.983 & 91 \\
M2 & 0.936 & 0.945 & 84 \\
M3 & 0.987 & 0.988 & 91 \\
M4 & 0.984 & 0.986 & 91 \\
B1 & 0.980 & 0.984 & 91 \\
B2 & 0.942 & 0.952 & 90 \\
\bottomrule
\end{tabular}
\end{table}

All OpenAlex differences in Table~\ref{tab:ci} except B2's lower bound ($-0.001$) exclude zero.
The noise floor is an order of magnitude smaller than the effects it is compared with: dropping
20\% of edges moves A1 by at most $0.014$ in $\rho$, whereas the smallest paired difference we
report in magnitude, $0.016$ for M3, and the M1 difference of $0.083$, are several times larger.
The observed between-metric spread is therefore not an artefact of missing edges.

\subsection{Semantic Scholar cross-check}

\begin{table}[h]
\caption{Semantic Scholar (13,548 edges, 1,456 authors).}
\label{tab:s2}
\centering
\footnotesize
\setlength{\tabcolsep}{3.5pt}
\begin{tabular}{lrrrrr}
\toprule
\textbf{Metric} & $\rho_{\text{fut}}$ & $\text{corr}_{\text{yr}}$ & \textbf{Top-10} & \textbf{Top-100} & \textbf{Ties} \\
\midrule
A1 & 0.502 & $-$0.002 & 24.2 & 9.1  & 82\% \\
M1 & 0.456 & $+$0.010 & 2.8  & 5.0  & 72\% \\
M2 & 0.126 & $+$0.090 & 4.1  & 4.5  & 99\% \\
M3 & 0.501 & $-$0.008 & 21.4 & 9.4  & 89\% \\
M4 & 0.126 & $-$0.003 & 6.2  & 5.5  & 67\% \\
B1 & 0.457 & $+$0.089 & 7.7  & 4.8  & 85\% \\
B2 & 0.223 & $+$0.043 & 22.5 & 7.0  & 100\% \\
\bottomrule
\end{tabular}
\end{table}

The \emph{order} of metrics by $\rho_{\text{fut}}$ is preserved only at the top: A1 $>$ M3 $>$
B2 $>$ M4 on OpenAlex becomes A1 $\approx$ M3 $>$ B1 $\approx$ M1 $>$ M2 $\approx$ M4 on Semantic
Scholar. A1 and M3 remain the two most accurate and M4 the least, but B2 and M1 swap places, and
M2 overtakes M4. The H1 result is reproduced in sign (M1's $\text{corr}_{\text{yr}}$ is less
negative, i.e.\ positive, against A1's $-0.002$), but both values lie within $\pm 0.01$ of zero,
so the magnitude of the OpenAlex move cannot be checked here. That is a limitation of the mirror,
not evidence against the OpenAlex result.

\subsection{Citation deflator (H5--H8)}

\begin{table}[h]
\caption{Deflator results. D1 = deflated total share, D2 = deflated per-paper share,
  B1$'$ = deflated per-paper citations.}
\label{tab:deflator}
\centering
\footnotesize
\setlength{\tabcolsep}{3.5pt}
\begin{tabular}{lrrrr}
\toprule
\textbf{Metric} & $\rho_{\text{fut}}$ & $\text{corr}_{\text{yr}}$ & $\rho$ vs A1 & \textbf{Overlap} \\
\midrule
\multicolumn{5}{l}{\textit{OpenAlex}} \\
A1  & 0.708 & $-$0.178 & 1.000 & 100 \\
M1  & 0.625 & $-$0.185 & 0.937 & 48 \\
B1  & 0.646 & $-$0.080 & 0.923 & 47 \\
D1  & 0.694 & $-$0.217 & 0.998 & 97 \\
D2  & 0.610 & $-$0.232 & 0.936 & 47 \\
B1$'$ & 0.634 & $-$0.132 & 0.929 & 45 \\
\midrule
\multicolumn{5}{l}{\textit{Semantic Scholar}} \\
A1  & 0.502 & $-$0.002 & 1.000 & 100 \\
M1  & 0.456 & $+$0.010 & 0.984 & 53 \\
B1  & 0.457 & $+$0.089 & 0.957 & 36 \\
D1  & 0.499 & $-$0.008 & 1.000 & 96 \\
D2  & 0.456 & $+$0.003 & 0.985 & 55 \\
B1$'$ & 0.457 & $+$0.073 & 0.964 & 38 \\
\bottomrule
\end{tabular}
\end{table}

\begin{table}[h]
\caption{Deflator curve on OpenAlex: outside-pool references per pool paper, and the resulting
  weight. Years before 2005 are excluded (truncated references).}
\label{tab:rt}
\centering
\footnotesize
\setlength{\tabcolsep}{4pt}
\begin{tabular}{lrr|lrr}
\toprule
\textbf{Year} & $R(t)$ & $D(t)$ & \textbf{Year} & $R(t)$ & $D(t)$ \\
\midrule
2005 & 18.43 & 1.742 & 2014 & 25.69 & 1.250 \\
2006 & 19.69 & 1.631 & 2015 & 28.03 & 1.145 \\
2007 & 20.61 & 1.558 & 2016 & 25.99 & 1.235 \\
2008 & 20.66 & 1.554 & 2017 & 27.57 & 1.165 \\
2009 & 19.70 & 1.630 & 2018 & 28.34 & 1.133 \\
2010 & 20.88 & 1.538 & 2019 & 25.94 & 1.238 \\
2011 & 20.61 & 1.558 & 2020 & 25.28 & 1.270 \\
2012 & 22.38 & 1.435 & 2021 & 32.11 & 1.000 \\
2013 & 23.37 & 1.374 & 2022 & 28.67 & 1.120 \\
\bottomrule
\end{tabular}
\end{table}

\paragraph{H5: failed, in the wrong direction.} $|\text{corr}_{\text{yr}}|$ rises from 0.178 to
0.217 on OpenAlex, and from 0.002 to 0.008 on Semantic Scholar. The predicted mechanism is
inverted. Citations received inside the scoring windows (born 2006--2014, cited by 2009--2017) are
discounted by $D(t) = 1.63$ in 2006, falling to $1.17$ in 2017 (Table~\ref{tab:rt}), while
citations received from 2018 to 2022 carry weights of $1.13$ to $1.27$. The discount is therefore
applied most heavily to exactly the later-born papers that the fixed window already puts at a
disadvantage, and the correlation moves away from zero.

\paragraph{H6: held.} $\rho_{\text{fut}} = 0.694$ against A1's 0.708: a loss of 0.014 versus
M1's 0.083, six times smaller, and the top-10 mean barely moves (91.6 against 96.5). On Semantic
Scholar the loss is 0.003.

\paragraph{H7: failed on reordering.} $\rho(\text{D1}, \text{A1}) = 0.998$ and top-100 overlap 97
of 100, against the predicted 0.80--0.95 and below 70. The top nine authors hold identical
positions and the remainder exchange adjacent places (for example de Rijke $10 \leftrightarrow 11$
with D\'iaz). Across all 1,430 authors, 1,256 change rank at all, but 79 of the 100 authors in A1's
top 100 stay inside it and the largest move is $70$ places (Barbosa, $718 \to 788$). $D(t)$ ranges
from 1.00 to 1.74 and is nearly flat after 2012, so within the scoring window it behaves almost as
a constant.

\paragraph{H8: failed as stated.} D1's $\text{corr}_{\text{yr}}$ is $-0.217$ against B1$'$'s
$-0.132$, so the field share keeps the correlation \emph{larger} in magnitude, not smaller. H8
asked which of the two is closer to zero; tested as stated, it fails. The pair also shows that D1
is the more accurate (0.694 against 0.634) with a four-times higher top-10 mean (91.6 against
22.0). We report H8 as a failed prediction, and do not claim from these numbers that the field
share adds information beyond deflation: on OpenAlex D1 and B1$'$ correlate at $0.929$ and share
45 of 100, and on Semantic Scholar at $0.964$ with 38 of 100.

\subsection{The birth-window confound}
\label{sec:confound}

The outcome counts citations received in a \emph{fixed calendar window} (2018--2022). A paper born
in $y$ collects scored citations in $[y, y+3]$, and the two windows are disjoint only when
$y + 3 < 2018$. For papers born after 2014 they overlap or coincide. Earlier-born papers therefore
not only had more time to accumulate citations before the outcome window; they also have a
guaranteed gap in which to be cited, while a 2014 paper's scored window (2014--2017) is followed
immediately by the outcome window (2018--2022). The mean birth year in our sample is 2010.7, and
A1's top 10 average 2010.4, so the leading group sits at the centre of the distribution, not at
either end.

Three checks support this reading and place limits on it. First, the sensitivity study gives
$\text{corr}_{\text{yr}} = -0.197, -0.189, -0.178, -0.163, -0.166$ for $\Delta = 1, 2, 3, 5, 8$:
widening the scoring window narrows the correlation without removing it. Second, widening the
author cutoff from 2014 to 2016, which admits later-born papers, moves it from $-0.178$ to
$-0.130$. Third, the top-10 mean birth year (2010.4) lies \emph{below} the population mean
(2010.7); the correlation is carried by the body of the distribution and not by the leaders. We
therefore treat $\text{corr}_{\text{yr}}$ as a diagnostic of the scoring protocol, not a
headline bias measure, and we do not quantify how much of it is mechanical. A protocol in which
every author is scored at the same number of years since first publication would separate the
mechanical component from any behavioural one; it is the second item of future work.

\section{Discussion}

\paragraph{What the comparison establishes.} At this scale the standard corrections are
separable. Per-paper averaging is a real and costly change: it reorders half the top 100 and gives
up $0.083$ of rank accuracy. Deflation is not: it costs $0.014$, leaves the top nine authors in
place, and is therefore a reporting correction rather than a re-ranking. Averaging and the field
share are close but not redundant ($\rho = 0.937$, 48 of 100 shared), and the field share and the
deflator are more nearly redundant still ($\rho = 0.998$).

\paragraph{What it does not establish.} The age-related quantity that motivated H1 and H5 is
confounded with the calendar structure of the outcome (Section~\ref{sec:confound}). Neither the
M1 result nor the D1 result can be read as ``normalisation does not reduce age bias'': the
confound is present in A1 and in every variant, and it is of the same size as the effects we
observe. The literature separates a pure age effect from productivity and collaboration
\citep{huang2012academicage}, reports academic-age normalisation as an operational remedy
\citep{anonymous2026platform}, and documents citation inflation \citep{petersen2019inflation}; our design cannot
adjudicate between these accounts for this cohort.

\paragraph{Practical recommendation.} Report the field-share sum as the headline, report the
deflated variant alongside it when absolute counts across eras are compared, and report the
per-paper average as a separate question rather than a replacement. Say in the protocol what the
outcome window is, because the result depends on it.

\subsection{Limitations}

\begin{enumerate}
  \item \textbf{One field, one source, one window.} Five IR venues, one birth window
    (2006--2014), one 5-year outcome. The correlations are specific to the overlap structure of
    these windows.
  \item \textbf{The birth-window confound.} The outcome window is shared by all authors, so the
    relation between score and birth year mixes a mechanical and a behavioural component that
    this design cannot separate.
  \item \textbf{Deflator supply.} $R(t)$ is measured on outside-pool references, so it
    under-represents fields that IR stopped citing. Its shape is what the deflator uses, and a
    field-wide reference population would be preferable; that pull was not available, so the
    per-paper and per-distinct-work variants ($D_{\text{naive}}$, $D_{\text{distinct}}$) remain
    cross-checks rather than alternatives reported in the tables.
  \item \textbf{Thin mirror.} Semantic Scholar is thin enough that its $\text{corr}_{\text{yr}}$
    values lie within $\pm0.01$ of zero; it tests directions, not magnitudes.
  \item \textbf{Tie-dominated indices.} M2 and B2 are 82--100\% ties, so their rank correlations
    are computed on a heavily tied vector. No tie-breaking alternative was tested.
  \item \textbf{Single setting for shrinkage.} Only $k=3$ was tested for M4.
  \item \textbf{No external benchmark.} No author-level IR ground truth is public; the temporal
    holdout is the only reproducible target available.
\end{enumerate}

\subsection{Future work}

The evidence points at the protocol, not at further aggregation schemes.
\begin{enumerate}
  \item \textbf{Percentile indicators} \citep{waltman2016review, waltman2013percentile}: rank each paper in its
    cohort and average per author. This is the only standard family not yet built.
  \item \textbf{Fixed-age scoring}: score every author at the same number of years since first
    publication, which removes the window-overlap confound by construction.
  \item \textbf{Position-weighted credit} \citep{vavrycuk2018influential} and \textbf{PageRank}
    \citep{ding2011pagerank} scaled to this graph.
\end{enumerate}

\section{Conclusion}

We compared seven author-level metrics and three deflator variants against a held-out
future-citation outcome on a closed IR graph, with hypotheses fixed before the numbers. Per-paper
averaging costs $0.083$ in rank accuracy and does not turn into a freshness ranking. Citation
deflation costs $0.014$ and changes almost nothing in the top 100. The field share survives both.
The quantity these corrections were introduced to fix cannot be measured fairly with a shared
outcome window, because paper age and outcome visibility are mechanically entangled. The next
step is a protocol change, not another metric.

\medskip
\bibliographystyle{plainnat}
\bibliography{neurips_biblio}

\end{document}
